% !TEX program = pdflatex
\documentclass[12pt,a4paper]{article}

% Encoding and language
\usepackage[T1]{fontenc}
\usepackage[utf8]{inputenc}
\usepackage[english]{babel}

% Bibliography
\usepackage[round,authoryear]{natbib}

% URL handling
\usepackage{xurl}

% Typography and layout
\usepackage{lmodern,microtype,setspace,geometry}
\usepackage{graphicx,caption,float,placeins}
\usepackage{fancyhdr,hyperref}
\usepackage{textcomp,eurosym,newunicodechar}
\usepackage{titlesec}

% Page layout -- double spacing, min 2.5 cm margins
\geometry{margin=2.5cm}
\doublespacing

% Math
\usepackage{amsmath,amssymb}

% Tables
\usepackage{booktabs,longtable,array,tabularx,threeparttable,makecell,adjustbox,rotating}

% Unicode characters used in source data
\newunicodechar{α}{\ensuremath{\alpha}}
\newunicodechar{β}{\ensuremath{\beta}}
\newunicodechar{γ}{\ensuremath{\gamma}}
\newunicodechar{≤}{\ensuremath{\leq}}
\newunicodechar{≥}{\ensuremath{\geq}}
\newunicodechar{≈}{\ensuremath{\approx}}
\newunicodechar{−}{-}

% Hyperref
\hypersetup{hidelinks,breaklinks=true}

% Paragraph formatting
\setlength{\parindent}{1.25em}
\setlength{\parskip}{0.25em}
\setlength{\textfloatsep}{16pt plus 2pt minus 2pt}
\setlength{\intextsep}{14pt plus 2pt minus 2pt}

% Caption setup
\captionsetup{font=small,labelfont=bf,justification=justified,singlelinecheck=false}

% Header/footer
\pagestyle{fancy}
\fancyhf{}
\fancyfoot[C]{\thepage}
\renewcommand{\headrulewidth}{0pt}

% Section formatting -- bold titles with Arabic numerals
\titleformat{\section}{\normalfont\bfseries\large}{\thesection.}{0.6em}{}
\titlespacing*{\section}{0pt}{2.0ex plus .6ex minus .2ex}{0.8ex plus .2ex}

\titleformat{\subsection}{\normalfont\bfseries\normalsize}{\thesubsection}{0.6em}{}
\titlespacing*{\subsection}{0pt}{1.5ex plus .5ex minus .2ex}{0.6ex plus .2ex}

\titleformat{\subsubsection}{\normalfont\bfseries\small}{\thesubsubsection}{0.6em}{}

\setcounter{secnumdepth}{3}

% -----------------------------------------------------------------------
\begin{document}
% -----------------------------------------------------------------------

% =======================================================================
% FIRST PAGE: title, abstract, keywords, JEL classification
% =======================================================================
\thispagestyle{fancy}

% TODO: Add author affiliation block here
\begin{center}
{\large\bfseries Reproeconomics: The Economic Impact of State-Funded
Oocyte Cryopreservation}
\end{center}

\vspace{0.8em}

% TODO: Add acknowledgements if required

\noindent\textbf{Abstract}

\noindent Spain has experienced persistent fertility decline and
pronounced childbearing postponement. Recent political proposals to
publicly fund elective oocyte cryopreservation (OC) raise questions
about fiscal sustainability and cost-effectiveness. Drawing on Spanish
demographic data from 1976--2022, I employ an Interrupted Time Series
design to calibrate a demographic additionality parameter, embedded in
a decision-analytic cost framework incorporating utilisation behaviour,
age-dependent success probabilities, and generational accounting
returns. The marginal cost per additional live birth substantially
exceeds any plausible fiscal return under all empirically plausible
parameter combinations. Publicly funded OC cannot be justified on
pro-natalist or fiscal grounds; the relevant argument is redistribution
and access.

\vspace{0.5em}

\noindent\textbf{Keywords:} Elective oocyte cryopreservation;
Cost-effectiveness; Fertility postponement; Health policy; Fiscal
sustainability

\vspace{0.5em}

% TODO: Confirm or add JEL classification codes
\noindent\textbf{JEL Classification:} H51, I18, J13, H55

% =======================================================================
% MAIN TEXT begins on new page
% =======================================================================
\newpage

% -----------------------------------------------------------------------
\section{Introduction}\label{sec:introduction}
% -----------------------------------------------------------------------

Two-thirds of the world's population today resides in nations where
the replacement rate of 2.1, the conventional estimate of what is
required to maintain a steady population, falls below
\citep{worldbank2025}. Spain's case is particularly acute. According
to the most recent data from Spain's National Statistics Institute,
318,005 births occurred in 2024, the lowest number since the
statistics series started in 1941 and a decline of almost 25\% over
the previous ten years \citep{ine2025mnp,funcas2024}. Women gave birth
to their first child at an average age of 32.6 years, compared to 25.6
in 1980 and 33.2 if only Spanish women are analysed
\citep{funcas2024}.

Against this demographic backdrop, assisted reproductive technologies
(ART) have increasingly been perceived as a means to mitigate the
biological constraints associated with delayed childbearing. The use of
elective oocyte cryopreservation (OC) has expanded rapidly in Spain
over the last decade: according to data from the National Activity
Registry of the Spanish Fertility Society \citep{sef2024registro}, the
number of women undergoing egg freezing increased almost eighteenfold
between 2011 and 2021. Spain records the highest share of births to
mothers aged 40 and over in the European Union, accounting for 10.7\%
of all live births in 2021 \citep{blas2025}.

The Popular Party's policy statement, approved at its July congress,
includes for the first time public subsidies for egg freezing for women
who wish to delay motherhood \citep{blas2025,partidopopular2025}.
Until now, Social Security has only covered this service in cases of
illness that damages fertility; women have had to resort to private
clinics for elective purposes. The project is modelled after the
Galician government's programme, which provides publicly funded OC for
women aged 30--35 under reproductive health criteria
\citep{xunta2024}. Public funding for elective OC remains the
exception across European health systems, where coverage is typically
limited to medical indications such as fertility preservation prior to
gonadotoxic therapy \citep{shenfield2017,gambadauro2023}. Existing studies assess OC
either from a clinical perspective or under individual-level economic
models, but evidence remains scarce on its cost-effectiveness and
fiscal sustainability when implemented as a publicly funded policy and
evaluated at the population level.

This paper has three objectives. First, conduct a full economic
evaluation of the proposed programme: constructing an incremental
cost-effectiveness ratio (ICER) relative to the status quo of no
public subsidy, estimating the annual budget impact under plausible
uptake scenarios, and analysing the distributional consequences of
alternative public-funding designs using concentration-curve
methodology. Second, evaluate whether such a policy can be
self-financing in the long run, in the sense that additional births
generated by the intervention yield an increase in aggregate output at
least comparable to programme costs. Third, determine whether publicly
funded access to elective OC is likely to increase fertility outcomes
in the medium term---a question whose answer bounds the demographic
additionality parameter central to the economic model.

The remainder of the paper is structured as follows.
Section~\ref{sec:literature} reviews the relevant literature.
Section~\ref{sec:methodology} outlines the methodological framework
and identification strategy. Section~\ref{sec:demographic} presents
the fertility dynamics results that calibrate the demographic
additionality parameter. Section~\ref{sec:economic} develops the
economic evaluation. Section~\ref{sec:discussion} discusses the
findings and their policy implications. Appendix~\ref{app:data}
provides data description and full regression output; Appendix~\ref{app:decisiontree}
presents a decision-tree representation of the cost model (see tables
and figures at the end of the manuscript).

% -----------------------------------------------------------------------
\section{Literature Review}\label{sec:literature}
% -----------------------------------------------------------------------

The literature converges on two findings central to the policy
question examined here. Access to ART is extremely uneven and strongly
correlated with socioeconomic position.
\citet{alonpinilla2021} demonstrate, using microdata from the Spanish
\emph{Encuesta de Fecundidad}, that access to ART and the choice
between public and private provision are significantly influenced by
income and education: treatment in private clinics is substantially
more common for women with higher incomes and levels of education,
while women from poorer socioeconomic backgrounds are more likely to be
denied access even when they report reproductive issues (2021:9).
Results also differ by sector, with women treated solely in the private
sector displaying a greater cumulative pregnancy incidence than those
treated in the public system (2021:15). Similar dynamics are documented
at the European level by \citet{seiz2023}, who find significant
cross-national disparities consistently and negatively affecting women
from lower socioeconomic backgrounds.

The demand for and attitudes towards elective OC help to clarify the
access dimension. The most frequently stated causes of non-medical OC,
according to a systematic review by \citet{kynigopoulou2024}, are the
lack of a suitable partner, investments in education and profession,
and general economic uncertainties. OC is widely viewed as
reproductive insurance---a purchase that reduces anxiety but is rarely
exercised. \citet{gambadauro2023}, drawing on survey data from Swedish
university students, find that OC is broadly accepted as a technology,
but that support for public funding is significantly lower when OC is
presented as age-related rather than medically required; there are no
significant differences in willingness to pay between scenarios,
implying that debate is more about who should pay than about the
perceived usefulness of the technology.

From an economic standpoint, the effectiveness and therefore the
efficiency of elective OC depends critically on two parameters: age at
freezing and subsequent utilisation. \citet{chronopoulou2021} show
that live birth rates per vitrified oocyte decline sharply with age,
making earlier freezing substantially more effective from a clinical
perspective. However, clinical potential does not always convert into
actual results. \citet{yang2022}, in a retrospective cohort of 2,038
women followed over seven years, found that only 8.4\% returned to use
their vitrified oocytes, yielding a population-level birth rate of just
3.0\% relative to all women who froze. In a Spanish private clinic
sample of 1,241 women, \citet{cobo2018} reported a 12.1\% return rate
over a median follow-up of 4.2 years; rates are as low as 5\% in
United Kingdom private settings \citep{trokoudes2011}. The cost per
live delivery rises significantly when expenditures are averaged among
all freezing women rather than only those who eventually use OC
\citep{yang2022}. Because utilisation is revealed-preference behaviour
that a price subsidy will only partially shift, programme cost must be
spread over all recipients, not only those who ultimately attempt a
pregnancy.

The formal cost-effectiveness literature reaches qualified conclusions.
\citet{loendersloot2011}, in a Markov model contrasting elective OC at
age 35 with no freezing followed by subsequent IVF, find that OC is
cost-effective only under stringent assumptions, requiring favourable
success probabilities and high utilisation rates simultaneously.
\citet{bakkensen2022}, drawing on US SART-CORS data, find that OC at a
younger age can lower the expected cost per live delivery and raise the
likelihood of reaching the intended family size among women who delay
childbearing, but sensitivity analyses reveal these findings depend
substantially on assumptions about freezing age and future use.
Positive outcomes at the individual level do not always translate to a
population-wide, publicly financed solution.

A particularly relevant body of literature takes a public finance
perspective. \citet{matorras2016} estimate the lifetime net fiscal
contribution (NFC) of individuals born through assisted reproduction in
Spain using a generational accounting methodology. In the base case,
the NFC of an IVF-conceived individual is \euro{}66,709 (at 2006
prices), implying a return of \euro{}15.98 per euro of public
investment; the discounted value of future taxes and social
contributions surpasses public spending despite significant upfront
medical expenditures (2016:118). This figure serves as the principal
benchmark against which the programme cost model in
Section~\ref{sec:economic} is calibrated. Critically, however,
\citet{matorras2016} model clinical IVF for infertile couples; the
parameters governing cost and success for elective OC are materially
different: the upfront vitrification cycle adds a cost category absent
from IVF, and the low utilisation rate means the effective cost per
birth is spread over a much larger denominator. International evidence
similarly shows positive lifetime fiscal contributions for ART-conceived
individuals in most countries, with the notable exception of the
Netherlands \citep{moolenaar2014}.

What is largely unexplored is whether publicly funded elective OC can
be justified once behavioural responses, selection into utilisation,
and long-term fiscal effects are considered jointly. The present paper
fills this gap by embedding empirically calibrated utilisation and
additionality parameters within a complete generational accounting
framework adapted specifically for elective OC.

% -----------------------------------------------------------------------
\section{Methodology}\label{sec:methodology}
% -----------------------------------------------------------------------

\subsection{Analytical framework}\label{sec:framework}

The analysis comprises two integrated components. The first employs an
Interrupted Time Series (ITS) design based on annual Spanish
demographic data (1976--2022) to characterise the structural change in
fertility dynamics at an endogenously determined breakpoint
$T_0 = 2010$. This component is explicitly a calibration exercise
rather than a causal evaluation, since no real OC programme existed
during the estimation window: the ITS identifies structural dynamics in
the historical fertility series, not the effects of an implemented
policy. The demographic findings from this component bound the
demographic additionality parameter $\alpha$, which is defined as the
share of OC-assisted births that would not have occurred by any other
means. The second component embeds these demographic findings in a
decision-analytic cost framework to estimate an ICER relative to the
status quo, assess fiscal self-sustainability, quantify budget impact
under alternative uptake scenarios, and evaluate the distributional
consequences of different funding designs.

\subsection{Cost parameters and assumptions}\label{sec:assumptions}

OC is best conceptualised as a sequence of clinically established yet
economically heterogeneous stages, as reliance on average treatment
costs risks obscuring substantial variation in marginal resource use.
The process begins with a gynaecologic diagnostic and eligibility
assessment. This is followed by controlled ovarian stimulation, during
which exogenous gonadotropins are administered over approximately
10--14 days; this phase typically constitutes the most
resource-intensive component. Once adequate follicular maturation is
achieved, oocyte retrieval is performed via ultrasound-guided
aspiration under light anaesthesia. Retrieved oocytes are subsequently
vitrified for long-term storage in liquid nitrogen \citep{pai2021}.

The official cost for assisted reproduction in the Spanish public
health system is not officially documented. Cost inputs are drawn from
published estimates of direct healthcare costs for IVF and artificial
insemination (AI) cycles in two Spanish institutions
\citep{matorrasweinig2011,matorras2001,navarro2006}. Based on an
average success rate of 28.7\% \citep{sef2001}, the expected investment
per live birth is \euro{}4,173 for IVF and \euro{}3,629 for AI (2006
prices), with delivery costs excluded as they are presumed constant
between assisted and naturally conceived children. Three caveats
qualify the direct application of these figures to elective OC. First,
all \citet{matorras2016} values are in 2006 prices; Spanish healthcare
input costs rose by about 37\% between 2006 and 2023 (INE IPCSNA), so
nominal comparisons understate present costs. Second, the 28.7\%
success rate reflects clinical IVF for infertility and is not directly
transferable to elective OC, since age at retrieval, age at thaw, and
number of vitrified oocytes interrelate non-linearly
\citep{cobo2016,chronopoulou2021}. Third, elective OC adds a cost
element absent from clinical IVF: the vitrification cycle itself
(stimulation, retrieval, laboratory processing, multi-year
cryostorage), approximated at \euro{}3,000--5,000 at SEF 2023 prices.
This cost is incurred regardless of whether the woman ever uses her
oocytes, making it a fixed programme cost that must be included in any
complete budget-impact analysis.

\subsection{The ITS design as a calibration tool}\label{sec:its}

When randomised assignment is not feasible and no reliable control
group exists, an ITS design is a segmented-trend regression framework
that can support causal inference when a genuine policy intervention is
externally assigned and the parallel-trends assumption holds
\citep{stockwatson2020}. In the present study, however, no actual
policy intervention took place during the estimation window. The ITS is
therefore applied in a calibration mode: the structural break
identified at $T_0 = 2010$ serves as a hypothetical anchor for the
counterfactual, and the estimated post-break coefficients characterise
the fertility dynamics that a real OC programme implemented at that
date would have needed to replicate.

Let $Y_t$ denote an aggregate outcome observed at time
$t = 1,\ldots,T$, and let $T_0$ be the structural breakpoint. Define
the binary break indicator

\begin{equation}
D_t =
\begin{cases}
0 & \text{if } t < T_0,\\
1 & \text{if } t \geq T_0.
\end{cases}
\label{eq:Dt}
\end{equation}

The break indicator $D_t$ does not represent an observed policy reform;
it marks the structural breakpoint at $T_0 = 2010$. Employing the
notation and arrangement of \citet{stockwatson2020} for policy analysis
with time trends, the ITS regression model is

\begin{equation}
Y_t = \beta_0 + \beta_1 t + \beta_2 D_t + \beta_3 (t - T_0) D_t + u_t
\label{eq:its-model}
\end{equation}

\noindent where $u_t$ is an unobserved error term. $\beta_1$ shows the
linear pre-break time trend, $\beta_2$ reflects the instantaneous level
shift at the structural break, and $\beta_3$ indicates the change in
the slope of the post-break trend. As emphasised by
\citet{stockwatson2020}, serial correlation does not bias OLS
coefficient estimates but invalidates conventional standard errors;
inference therefore relies on heteroskedasticity- and
autocorrelation-consistent (HAC) standard errors throughout. All
post-break parameters must be interpreted as identifying the structural
change in fertility at $T_0$, not determining the causal impact of a
programme.\footnote{The logit specification described in
Appendix~\ref{app:data} complements the linear ITS as a qualitative
robustness check by evaluating the probability of observing
high-fertility outcomes before and after the break. Owing to
quasi-complete separation arising from the endogenous construction of
its dependent variable, the logit results should not be interpreted as
an independent source of identification.}

\subsection{Intervention year selection}\label{sec:intervention-year}

In any ITS design, the choice of the intervention year is crucial as
it controls the temporal partition between pre-intervention trends and
post-intervention dynamics. The intervention year is $T_0 = 2010$,
derived from a formal Euclidean distance criterion. The 2020
observation is treated as the terminal observation and compared with
each preceding year in the dataset; the distance-minimisation criterion
identifies the year in which Spain's fertility structure most closely
resembled its 2020 configuration, thereby marking the beginning of the
structural regime characterising the post-break period. The pre-trend
test in Figure~\ref{fig:pretrend} confirms close alignment between the
observed trajectory and the estimated pre-break linear trend using only
pre-2010 data, supporting the counterfactual assumption.

The break-date sensitivity analysis in Figure~\ref{fig:break_sensitivity}
provides additional justification for $T_0 = 2010$. Specifications
allocating the structural break to earlier candidate years produce
smaller, less statistically significant level-shift coefficients. At
$T_0 = 2010$, the estimated level-shift coefficient is the largest in
absolute magnitude, negative, and highly statistically significant
across all candidate years.\footnote{This does not establish 2010 as
the causal date of any policy; it confirms that fertility data show
their sharpest discontinuity at that point in the sample period.}
$T_0 = 2010$ provides a pre-break window of 34 observations
(1976--2009) and a post-break window of 13 observations (2010--2022),
sufficient to identify the baseline trend reliably and estimate slope
corrections. This framing governs every interpretative claim in
Sections~\ref{sec:demographic} and~\ref{sec:economic}: post-break
dynamics reflect the fertility dynamics linked to the historically
recognised structural break, not the causal impact of an established
OC programme.

% -----------------------------------------------------------------------
\section{Fertility Dynamics and Additionality Calibration}\label{sec:demographic}
% -----------------------------------------------------------------------

The ITS results reported here capture the deviation of observed
fertility from the pre-2010 counterfactual trajectory and are used to
calibrate the demographic additionality parameter $\alpha$ entering the
economic model. The central diagnostic question is whether the
post-break dynamics exhibit the age-specific recovery that a properly
functioning OC programme would predict: partial recuperation of
fertility at ages 35--39 (oocytes used after a delay) and continued
fertility at ages 40--44 (successful use within the target age window).
Whether delayed childbearing translates into higher completed fertility
depends on whether later-life biological constraints outweigh the gains
from postponement.

For women aged 35--39, the post-break gap between fitted and
counterfactual age-specific fertility rates (ASFRs) is negative in the
years immediately following 2010, indicating that observed fertility
falls below the level predicted by continuation of the pre-break trend.
Over time, however, the magnitude of this gap progressively declines
and in later post-break years approaches zero and becomes slightly
positive, a pattern consistent with partial delayed-childbearing
recuperation at this age group.

For women aged 40--44, the pattern is qualitatively different. The
post-break gap remains predominantly negative throughout the entire
post-2010 period, with no sustained convergence toward the
counterfactual trajectory (Figure~\ref{fig:asfr_gap}). This
asymmetry---partial recovery at ages 35--39 but a permanent negative
deviation at ages 40--44---is the central demographic finding used to
calibrate $\alpha$.

This asymmetry has a specific relevance for the economic evaluation.
If an OC programme targets women aged 30--35, their oocytes would
typically be used at ages 38--45, exactly the age window where the
ASFR remains below its counterfactual. This suggests that even women
who successfully use their oocytes would be operating within conditions
of structural fertility suppression, reducing the likelihood that the
resulting birth is genuinely additional. The aggregate evidence further
indicates that the structural break at $T_0 = 2010$ is associated with
a large and statistically significant negative level shift in adjusted
total fertility (adjTFR) with no subsequent convergence toward the
pre-break counterfactual; the post-break deviation remains negative
over the entire post-2010 period. This pattern does not conform with
complete intertemporal substitution and suggests that delayed
childbearing does not lead to improved total fertility levels at the
aggregate level. The demographic additionality parameter $\alpha$
should therefore be calibrated toward its lower range
($\alpha \in [0.3,\,0.5]$), which mechanically raises the estimated
cost per additional birth. Full regression results, including the
linear ITS estimates with macroeconomic controls, are reported in
Table~\ref{tab:its-linear} in Appendix~\ref{app:data} (see tables at
the end of the manuscript).

% -----------------------------------------------------------------------
\section{Economic Evaluation}\label{sec:economic}
% -----------------------------------------------------------------------

\subsection{Generational accounting benchmark}\label{sec:benchmark}

Infertility affects 15\% of Spain's population, approximately one in
seven couples of reproductive age \citep{matorrasweinig2011}. The
Spanish public health care system provides two approaches to assisted
reproduction: IVF and AI, with public financing covering a maximum of
three IVF cycles or four AI cycles. Eligibility is limited to females
between the ages of 18 and 40 who are classified as infertile, with
priority given to nulliparous infertile couples \citep{alberto2002}.
Despite this framework, demand outstrips supply: in 2010, about one
million couples requested treatment but just 22\% received at least one
cycle, 65\% of which was done in private institutions. Average waiting
times in the public system reached 339 days, creating both extra
demand and inequity in access. This perceived access gap is the
appropriate starting point for an equity argument in favour of publicly
funded elective OC: if clinical IVF is already publicly funded for
medical indications, rejecting elective OC on cost grounds while
tolerating a 14-fold socioeconomic gradient in private access
\citep{alonpinilla2021} presents a tension with the Spanish National
Health System (SNS) principle of universal coverage. Yet equity in
access does not justify public expenditure without evidence that the
programme accomplishes meaningful demographic outcomes at a reasonable
cost.

From a macro-fiscal perspective, previous research suggests that even
under optimistic estimates, a shrinking birth rate and population
ageing can produce fiscal imbalances mitigable only by cuts in
government spending or increases in taxation
\citep{imw2012,mess2006,nrc2012,perezcamarero2012}. Higher fertility
would contribute to a medium-term growth of the taxpaying population,
partially offsetting the fiscal pressures of ageing. Generational
accounting models estimate the lifetime NFC of a person born through
ART as the present value of all gross revenues received by the State
from the individual over the life cycle minus direct transfer
expenditures, with transfers and taxes apportioned by the risk of
falling into various socioeconomic classes in each year of life.

Generational accounting findings indicate that the lifetime present
value of taxes paid by a naturally conceived individual is
\euro{}370,482, substantially larger than transfers of \euro{}275,972,
yielding an NFC of \euro{}94,510. After treatment costs, the NFC of an
IVF-conceived individual is \euro{}66,709 and \euro{}67,253 for AI (at
2006 prices), implying returns of \euro{}15.98 and \euro{}18.53 per
euro invested, respectively \citep{matorras2016}. All cost data from
\citet{matorras2016} are denominated in 2006 euros; applying the INE
IPCSNA healthcare cost deflator of 37\% (2006--2023) yields an
inflation-adjusted NFC of approximately \euro{}91,391 and an adjusted
expected IVF investment per live birth of approximately \euro{}5,717 at
2023 prices.

The analytic step not undertaken by \citet{matorras2016} is to embed
the NFC estimate within a complete cost-effectiveness model that
incorporates: (i) the vitrification cost layer specific to elective OC;
(ii) the empirically observed utilisation rate, which determines how
many women must freeze oocytes per live birth; (iii) the age-stratified
success rate at thaw; and (iv) the demographic additionality factor
$\alpha$. Once these four adjustments are applied, the effective fiscal
return per euro of public OC expenditure falls considerably below the
\euro{}15.98 reported for clinical IVF and under most empirically
plausible parameter combinations falls below the fiscal break-even
point.

\subsection{Parameters: utilisation rates and success rates}\label{sec:parameters}

The percentage of women who actually use their vitrified oocytes is the
sole major driver of overall programme cost-effectiveness. Recorded
utilisation rates in published cohort studies are much lower than
programme designers tend to assume. In the retrospective cohort of
2,038 women spanning seven years, \citet{yang2022} found that only
8.4\% returned to use their oocytes, with a live birth rate of 3.0\%
relative to all women who froze. \citet{cobo2018} reported a 12.1\%
return rate in a Spanish private clinic sample of 1,241 women over a
median follow-up of 4.2 years. Rates are as low as 5\% in United
Kingdom private settings \citep{trokoudes2011}. These findings are
consistent with qualitative evidence in \citet{kynigopoulou2024}: OC
acts largely as reproductive insurance---a purchase that reduces
anxiety but is rarely exercised in practice. Such low utilisation is
not a market failure amenable to subsidy; it reflects the preferences
of users who value the option, not the outcome.

The 28.7\% average success rate for IVF cycles conducted for clinical
infertility \citep{sef2001} does not transfer directly to elective OC
for two reasons. First, elective OC is conducted most often between
ages 30 and 37, with oocytes typically not used until ages 38--44,
where biological age at retrieval and age at utilisation co-determine
success. Second, the number of mature oocytes vitrified per cycle
moderates success continuously: \citet{cobo2016} estimate that a woman
who vitrifies five oocytes has a cumulative live birth probability of
approximately 22\% (at age $<35$), compared to 72\% if fifteen or more
oocytes are obtained. This highlights that the frequently reported
average success rate substantially overestimates outcomes for women who
freeze fewer than ten oocytes or do so after age 37, the modal profile
of elective users. These age and quantity gradients underpin the
parameter ranges used in the sensitivity analysis below.

\subsection{Cost per additional live birth}\label{sec:cost-per-birth}

Let $C_u$ denote the cost of a vitrification cycle including five-year
storage (\euro{}5,000 at SEF 2023 prices), $C_T$ the cost of a
thaw-and-transfer IVF cycle (\euro{}2,500), $u$ the utilisation rate,
$s$ the success rate per thaw cycle, and $\alpha$ the demographic
additionality factor. The expected cost per additional live birth is

\begin{equation}
\text{Cost per birth} = \frac{C_u + u\,C_T}{u\,s\,\alpha}.
\label{eq:cost-per-birth}
\end{equation}

At the base-case parameters ($u = 0.08$, $s = 0.30$, $\alpha = 0.50$),
the cost per additional live birth is \euro{}433,333. Figure~\ref{fig:cost_sensitivity}
presents a bidimensional sensitivity analysis over the plausible ranges
of $u$ and $s$, with $\alpha$ held constant at 0.50. Costs range from
\euro{}137,500 (optimistic: $u = 0.20$, $s = 0.40$) to above
\euro{}1,000,000 (pessimistic: $u = 0.05$, $s = 0.20$). At no point
in the plausible parameter space does the cost per birth fall below
the \euro{}91,391 inflation-adjusted net fiscal contribution per
additional birth. In the empirical utilisation range (5--12\%), the
cost per additional birth lies between \euro{}433,333 and
\euro{}683,333, far exceeding any defensible willingness-to-pay
threshold for a publicly funded pro-natalist intervention.

The tornado chart in Figure~\ref{fig:tornado} illustrates that the
utilisation rate is by far the primary uncertainty driver: varying $u$
from 5\% to 20\% alone shifts the cost per birth by more than
\euro{}500,000, dwarfing the sensitivity to any other parameter.

\FloatBarrier

\subsection{Fiscal self-sustainability condition}\label{sec:fiscal}

The comparison of cost per additional birth to the
inflation-adjusted NFC from \citet{matorras2016} illustrates the
financial arithmetic of the programme. At the base-case parameters
($u = 0.08$, $s = 0.30$, $\alpha = 0.50$), the cost per additional
live birth is \euro{}433,333 and the gross fiscal return per euro of
public expenditure is approximately \euro{}0.21 (\euro{}91,391 /
\euro{}433,333), well below break-even and dramatically less than the
\euro{}15.98 return for clinical IVF. Under the pessimistic utilisation
scenario ($u = 0.05$, $s = 0.20$), costs per birth exceed
\euro{}1,025,000 and the fiscal return falls to \euro{}0.09 per euro
invested, implying a structural fiscal loss of more than \euro{}0.91
per euro spent.

The fiscal self-sustainability condition requires

\begin{equation}
u\,s\,\alpha \;\geq\; \frac{C_u + u\,C_T}{\mathrm{NFC}}.
\label{eq:fiscal-condition}
\end{equation}

Because $u$ appears on both sides, this condition is nonlinear and must
be evaluated numerically. At the inflation-adjusted NFC of
\euro{}91,391, holding $\alpha = 0.60$ and $s = 0.35$ at their upper
plausible bounds, the break-even utilisation rate is
$u \approx 0.30$---approximately 3.8 times the empirical average.
Under the more conservative 2006 NFC of \euro{}66,709, the required
utilisation rate rises to $u \approx 0.43$, more than five times the
observed average. The programme therefore fails the fiscal
self-sustainability condition on either price basis by a large margin
at any empirically plausible parameter combination.

\FloatBarrier

\subsection{Budget impact}\label{sec:budget}

It is estimated that Spain has approximately 1.8 million women aged
30--35 \citep{ine2023}. Figure~\ref{fig:budget_scenarios} illustrates
three uptake scenarios and the corresponding annual budgetary
commitment. At 5\% uptake, nearly 90,000 women would freeze annually;
at the stated vitrification cost of \euro{}5,000 per cycle, gross
vitrification expenditure alone amounts to \euro{}450M per year. Adding
expected thaw-and-transfer costs (approximately \euro{}18M, based on 8\%
utilisation and \euro{}2,500 per cycle) and administration (\euro{}2M),
annual programme costs total approximately \euro{}470M. Under a 10\%
uptake scenario, annual costs reach approximately \euro{}938M; at 20\%
uptake, they exceed \euro{}1.87Bn, equivalent to nearly 2.3\% of total
Spanish public health expenditure \citep{minsan2024}. The costs of
treating adverse events (Section~\ref{sec:adverse}) are excluded from
these estimates and would add a further 5--8\% to direct procedure
costs.

The scale of these budgetary commitments warrants explicit note. Spain
currently spends around \euro{}1.9Bn each year on family and child
policy \citep{oecd2024}. Comparative demographic evidence indicates
that affordable childcare, wage-replacement parental leave, and housing
support interventions tend to achieve stronger and more consistent
effects on completed fertility compared to ART-based programmes
\citep{gauthier2007,rindfuss2010}. Even the minimum (5\% uptake)
scenario would consume approximately \euro{}470M per year---equivalent
to roughly 25\% of Spain's entire annual family and child policy
budget---for a programme whose cost per additional birth exceeds the
per-birth net fiscal contribution under all empirically plausible
parameter values. This opportunity cost should be a key factor in any
cost-benefit consideration of the measure.

\subsection{Distributive analysis}\label{sec:distributive}

The distributive implications of public funding for elective OC depend
on two factors: the economic profile of existing private users and the
design of the public programme.
Figure~\ref{fig:quintiles} shows estimated OC utilisation rates by
income quintile in the current private market and under two public
funding designs, calibrated on the access gradients documented by
\citet{alonpinilla2021} and \citet{seiz2023}. The statistics expose a
sharply pro-rich profile in the current private market: the estimated
utilisation rate in the top income quintile is approximately fourteen
times that of the bottom quintile. A universal public subsidy reduces
this gradient but does not eliminate it, given that informational
barriers, geographic access, and cultural attitudes towards ART are
consistently associated with socioeconomic status. A means-tested
subsidy allocated to women in income quintiles Q1--Q3 achieves a near
neutral distribution of programme gains while substantially reducing
the fiscal burden.

Figure~\ref{fig:concentration} shows stylised concentration curves
\citep{kakwani1977} corresponding to each of the three designs. The
private market yields a concentration index of approximately 0.42,
strongly pro-rich, calibrated on Spanish income quintile shares (World
Bank WDI 2023) and the empirically documented Q5/Q1 access gradient of
14 \citep{alonpinilla2021}. Universal public funding reduces this
index to approximately 0.18; means-tested funding approaches
distributional neutrality (concentration index $\approx 0.05$). If the
declared goal of public OC funding is to improve equity of access to
reproductive technology, then a means-tested instrument is clearly
preferable to a universal subsidy on both efficiency and equity
grounds.

Three alternative institutional designs should also be assessed. First,
direct public delivery through existing ART units---rather than a
demand-side subsidy to private clinics---would integrate the programme
with the existing governance and quality control regime in the SNS;
however, with average waiting times of 339 days, expanding eligibility
to elective OC without a matching capacity expansion would likely
prolong waiting times for clinically indicated IVF. Second, a
progressive co-payment system (analogous to Spain's pharmaceutical
co-payment) could support partial cost recovery from higher-income
users while preserving access for lower-income users. Third, a
mandatory pre-treatment counselling session on age-specific success
probabilities and the realistic likelihood of utilisation should be
included in any design, given that prospective users consistently
overestimate both parameters in the literature
\citep{kynigopoulou2024,gambadauro2023}.

\subsection{Adverse event costs}\label{sec:adverse}

Costs associated with adverse events should be included in any
comprehensive cost estimate, though they are excluded from the
base-case analysis. Ovarian hyperstimulation syndrome (OHSS) is the
major risk of controlled ovarian stimulation: moderate OHSS occurs in
3--6\% of stimulation cycles and severe OHSS requiring hospitalisation
in 0.5--2\% \citep{luke2012,pai2021}. In the base scenario of 180,000
cycles per year (10\% uptake), severe OHSS would generate
approximately 900--3,600 hospitalisations per year, at an average cost
of \euro{}3,500--8,000 per episode based on \citet{minsan2024} tariffs,
implying an annual adverse event cost of \euro{}3M--29M above direct
procedure costs. Follicular puncture complications occur in
approximately 0.3\% of retrievals; at programme scale these are not
negligible. A publicly funded programme should incorporate these
externalities into its budget envelope from the outset.

The long-run safety profile of children born from vitrified oocytes
has not yet been fully quantified, since vitrification protocols were
standardised only between 2006 and 2008 and long-term follow-up data
are restricted to approximately 15 years. Present evidence does not
show higher rates of congenital anomalies or adverse developmental
outcomes compared to conventionally conceived children
\citep{cobo2022,noyes2009}. However, existing evidence is not
sufficiently powered to detect rare outcomes at the frequency that
would become epidemiologically visible if a large public programme
dramatically expanded the cohort of vitrification-conceived individuals.
A publicly funded programme of the scale outlined here would essentially
be a population-level experiment whose long-term safety data would not
be available for two decades---an irreversibility that any precautionary
policy analysis must consider.

\subsection{Provider incentives}\label{sec:providers}

The Spanish ART market is predominantly private: in 2021, approximately
78\% of all ART cycles were undertaken in private institutions
\citep{sef2024registro}. Private ART providers face a structural
financial incentive to expand the uptake of elective OC, which is a
high-margin, recurring-revenue service with relatively low marginal cost
per additional patient in facilities already staffed to provide ovarian
stimulation. The introduction of a public subsidy would expand the pool
of solvent demand by reducing price barriers, directly increasing
private-clinic revenues.

This creates a classic supplier-induced demand risk \citep{evans1974}.
\citet{deonandan2012} document that ART providers systematically
exaggerate the odds of successful treatment in counselling interactions,
and that the marketing of elective OC tends to conflate the probability
of having frozen oocytes with the probability of a live birth---a
distinction that is highly consequential in practice. A demand-side
subsidy channelled through private firms without price regulation,
compulsory outcome disclosure, or independent counselling obligations
would expose the public payer to moral hazard from both the supply side
(demand created by providers) and the demand side (women who freeze but
would never have used their oocytes, creating pure deadweight cost).
The regulatory conditions to minimise these risks---tariff regulation,
accreditation requirements, compulsory pre-treatment counselling by an
independent clinician, public reporting of clinic-level outcomes
disaggregated by age and number of oocytes---are absent from the
current Spanish regulatory environment and from the PP policy proposal
\citep{partidopopular2025}. Their absence is not a minor implementation
detail but a structural flaw that would drive programme costs upward
without improving outcomes.

\subsection{Incremental cost-effectiveness ratio}\label{sec:icer}

The analysis in the preceding sections is summarised by a single ICER,
computed relative to the comparator of no public subsidy (status quo
of private financing). The incremental cost is the net public
expenditure per eligible woman under the proposed design; the
incremental effectiveness is the expected additional live births
attributable to the programme, controlling for demographic
additionality. The base-case ICER for publicly funded elective OC is
\euro{}433,333 per additional live birth, with a sensitivity range of
\euro{}137,500 (optimistic) to above \euro{}1,000,000 (pessimistic).

There is no established willingness-to-pay threshold for live births as
a policy outcome in the Spanish health system, which operates on a
quality-adjusted life year (QALY) framework with an implicit threshold
of approximately \euro{}25,000--30,000 per QALY \citep{vallejotorres2018}.
Applying QALY-value terminology requires assumptions on the
quality-of-life benefit of an additional child to society and to
parents; under the assumption of 20 full-quality life years per
expected birth, the base-case cost of \euro{}433,333 per birth is
equivalent to \euro{}21,667 per QALY, which lies within the
conventional threshold. However, this calculation is dominated by the
contested NFC projection and by the assumption of 20 full-quality life
years---assumptions that are highly uncertain and sensitive to
long-run macroeconomic conditions. Even where the QALY cost falls
within the threshold range, publicly funded elective OC cannot be
deemed justified on fiscal self-sustainability grounds at any
empirically meaningful parameter value: the break-even utilisation
rate under optimistic clinical assumptions is approximately 3.8 times
the empirical average, and this condition is effectively unattainable
under current realised parameters.

% -----------------------------------------------------------------------
\section{Discussion}\label{sec:discussion}
% -----------------------------------------------------------------------

This paper assessed whether publicly funded elective OC could raise
medium-term fertility, satisfy a fiscal sustainability condition, and
represent a cost-effective use of public expenditure relative to
alternative policy instruments. The analysis addresses these objectives
through two integrated components: an ITS design based on Spanish
demographic data for 1976--2022 that characterises structural fertility
dynamics and bounds the degree of demographic additionality required
for programme feasibility, and a complete economic assessment
comprising a decision-analytic cost model, ICER computation,
sensitivity analysis, budget impact projections, and distributional
analysis.

Age-specific effects reveal a sharp asymmetry with no evidence of
complete compensation. Women aged 35--39 exhibited a post-break gap in
ASFRs relative to the counterfactual that gradually narrowed and became
slightly positive, indicating partial recovery consistent with delayed
childbearing. By contrast, fertility at ages 40--44 remains permanently
below its counterfactual trajectory after 2010. The absence of recovery
at the oldest reproductive age intervals implies that postponement is
only partially offset and that some fertility losses are not recovered---
a pattern inconsistent with complete demographic additionality from an
OC programme targeting the 30--35 age range. At the aggregate level,
the structural break at $T_0 = 2010$ yields a significant level shift
in adjTFR with no convergence toward the pre-break counterfactual,
indicating not a short-run fluctuation but a structural change in the
fertility regime.

The economic assessment qualifies and extends these demographic
findings in three significant respects. First, the cost-per-additional-birth
analysis indicates that publicly funded elective OC is unlikely to
become a self-sustaining financial process under empirically plausible
utilisation and success rates. In the base case ($u = 0.08$,
$s = 0.30$, $\alpha = 0.50$), the cost per extra live birth is
approximately \euro{}433,333, compared with the inflation-adjusted NFC
of an OC-conceived individual of \euro{}91,391 at 2023 prices (updated
from the \citet{matorras2016} base case of \euro{}66,709 at 2006 prices
using the INE IPCSNA deflator). The gross fiscal return per euro spent
is therefore approximately \euro{}0.21---well below break-even and
dramatically below the \euro{}15.98 return estimated for clinical IVF
by \citet{matorras2016}. Under pessimistic but empirically verifiable
utilisation rates ($u = 0.05$, $s = 0.20$), costs per birth exceed
\euro{}1,025,000 and the fiscal return falls to approximately
\euro{}0.09 per euro, implying a structural fiscal loss of more than
\euro{}0.91 per euro of public expenditure. Second, the cost of a
national programme would run to approximately \euro{}470M (5\% uptake)
to over \euro{}1.87Bn (20\% uptake) per year, with adverse event costs
not included in the direct procedure budget. Third, the distributional
analysis demonstrates that a universal subsidy is moderately pro-rich,
while a means-tested design targeting income quintiles Q1--Q3 achieves
near distributional neutrality at a lower fiscal cost.

Three policy implications emerge. First, the benefits of publicly
funded elective OC should not be justified on pro-natalist or fiscal
sustainability grounds; the ITS evidence and cost model together suggest
that such justification would require implausibly optimistic assumptions
about both demographic impact and utilisation. If the policy moves
forward, it should be articulated on autonomy and equity principles---
reducing financial obstacles to a reproductive option that lower-income
women currently cannot access---with full transparency regarding the
low probability of future use and the only moderate increase in live
birth probability. Second, any programme must include compulsory
pre-treatment counselling by an independent clinician, public reporting
of age- and oocyte-stratified success rates, and tariff regulation to
discourage supplier-induced demand in the largely private ART market.
Third, the \euro{}470M--1.87Bn annual cost of an OC subsidy programme
should be compared against evidence-based pro-natalist policies such as
affordable childcare and wage-replacement parental leave, which have
better-documented demographic effectiveness and distributionally
progressive effects \citep{gauthier2007,rindfuss2010}.

The findings suggest a near correspondence between
postponement-facilitated reproductive technologies and aggregate
fertility outcomes: policies that relax biological timing constraints
in the absence of changing underlying fertility preferences are
unlikely to counter structural low fertility or achieve equilibrium in
fiscal settings. The salient issue of publicly funded OC is therefore
redistribution and access in the interest of policy, not a
pro-natalist mechanism to be deployed.

% =======================================================================
% BIBLIOGRAPHICAL REFERENCES
% =======================================================================
\clearpage
\begingroup
\small
\sloppy
\section*{Bibliographical references}
\addcontentsline{toc}{section}{Bibliographical references}

\begin{thebibliography}{99}

\bibitem[Alberto et~al.(2002)]{alberto2002}
Alberto, A., et~al. (2002). Guidelines for the use of resources
belonging to the Spanish National Health System in assisted reproductive
techniques. \emph{Revista Iberoamericana de Fertilidad}.

\bibitem[Alon and Pinilla(2021)]{alonpinilla2021}
Alon, I., \& Pinilla, J. (2021). Assisted reproduction in Spain,
outcome and socioeconomic determinants of access. \emph{International
Journal for Equity in Health}, 20, 156. DOI:
\nolinkurl{10.1186/s12939-021-01438-x}.

\bibitem[Bakkensen et~al.(2022)]{bakkensen2022}
Bakkensen, J. B., Flannagan, K. S., Mumford, S. L., Hutchinson, A. P.,
Cheung, E. O., Moreno, P. I., \ldots{} Goldman, K. N. (2022). A SART
data cost-effectiveness analysis of planned oocyte cryopreservation
versus in vitro fertilization with preimplantation genetic testing for
aneuploidy considering ideal family size. \emph{Fertility and Sterility},
118(5), 936--946.

\bibitem[Blas(2025)]{blas2025}
Blas, E. G. (2025, June 16). Feij\'oo hace un gu\~no progresista y
plantea que el Estado pague la congelaci\'on de \'ovulos. \emph{El Pa\'is}.

\bibitem[Chronopoulou et~al.(2021)]{chronopoulou2021}
Chronopoulou, E., Raperport, C., Sfakianakis, A., Srivastava, G., \&
Homburg, R. (2021). Elective oocyte cryopreservation for age-related
fertility decline. \emph{Journal of Assisted Reproduction and Genetics},
38, 1177--1186. DOI: \nolinkurl{10.1007/s10815-021-02072-w}.

\bibitem[Cobo et~al.(2016)]{cobo2016}
Cobo, A., Garc\'ia-Velasco, J. A., Coello, A., Domingo, J., Pellicer,
A., \& Remoh\'i, J. (2016). Oocyte vitrification as an efficient option
for elective fertility preservation. \emph{Fertility and Sterility},
105(3), 755--764.e8. DOI:
\nolinkurl{10.1016/j.fertnstert.2015.12.044}.

\bibitem[Cobo et~al.(2018)]{cobo2018}
Cobo, A., Garc\'ia-Velasco, J. A., Domingo, J., Pellicer, A., \&
Remoh\'i, J. (2018). Elective and onco-fertility preservation: factors
related to IVF outcomes. \emph{Human Reproduction}, 33(12), 2222--2231.
DOI: \nolinkurl{10.1093/humrep/dey321}.

\bibitem[Cobo et~al.(2022)]{cobo2022}
Cobo, A., Garc\'ia-Velasco, J. A., Remoh\'i, J., \& Pellicer, A.
(2022). Obstetric and perinatal outcome of babies born from vitrified
oocytes. \emph{Fertility and Sterility}, 117(1), 180--190.

\bibitem[Deonandan et~al.(2012)]{deonandan2012}
Deonandan, R., Green, S., \& van Beinum, A. (2012). Ethical concerns
for maternal surrogacy and reproductive tourism. \emph{Journal of
Medical Ethics}, 38(12), 742--745. DOI:
\nolinkurl{10.1136/medethics-2012-100551}.

\bibitem[Evans(1974)]{evans1974}
Evans, R. G. (1974). Supplier-induced demand: Some empirical evidence
and implications. In M. Perlman (Ed.), \emph{The Economics of Health
and Medical Care} (pp.\ 162--173). Macmillan.

\bibitem[Funcas(2024)]{funcas2024}
Fundaci\'on de las Cajas de Ahorros (Funcas). (2024). \emph{El
n\'umero de hijos por mujer en Espa\~na cay\'o en 2023 al m\'inimo
hist\'orico: 1,12}. Madrid. URL:
\url{https://www.funcas.es/prensa/el-numero-de-hijos-por-mujer-en-espana-cayo-en-2023-al-minimo-historico-112/}.

\bibitem[Gambadauro et~al.(2023)]{gambadauro2023}
Gambadauro, P., Br\"ann, E., \& Hadlaczky, G. (2023). Acceptance and
willingness-to-pay for oocyte cryopreservation in medical versus
age-related fertility preservation scenarios among Swedish female
university students. \emph{Scientific Reports}, 13, 5325.

\bibitem[Gauthier(2007)]{gauthier2007}
Gauthier, A. H. (2007). The impact of family policies on fertility in
industrialized countries: A review of the literature. \emph{Population
Research and Policy Review}, 26(3), 323--346. DOI:
\nolinkurl{10.1007/s11113-007-9033-x}.

\bibitem[Granger and Newbold(1974)]{grangernewbold1974}
Granger, C. W. J., \& Newbold, P. (1974). Spurious regressions in
econometrics. \emph{Journal of Econometrics}, 2(2), 111--120. DOI:
\nolinkurl{10.1016/0304-4076(74)90034-7}.

\bibitem[Instituto Max Weber(2012)]{imw2012}
Instituto Max Weber. (2012). \emph{Proyecciones demogr\'aficas y
sostenibilidad fiscal en Espa\~na}. Madrid.

\bibitem[INE(2023)]{ine2023}
Instituto Nacional de Estad\'istica. (2023). \emph{Population by age
and sex: women aged 30--35 in Spain}. Madrid.

\bibitem[INE(2025)]{ine2025mnp}
Instituto Nacional de Estad\'istica. (2025). \emph{Movimiento Natural
de la Poblaci\'on 2024}. Madrid. URL:
\url{https://www.ine.es/dyngs/Prensa/MNP2024.htm}.

\bibitem[Kakwani(1977)]{kakwani1977}
Kakwani, N. C. (1977). Measurement of tax progressivity: An
international comparison. \emph{The Economic Journal}, 87(345), 71--80.
DOI: \nolinkurl{10.2307/2231833}.

\bibitem[Kynigopoulou et~al.(2024)]{kynigopoulou2024}
Kynigopoulou, S., Matsas, A., Tsarna, E., Christopoulou, S.,
Panagopoulos, P., Bakas, P., \& Christopoulos, P. (2024). Egg
cryopreservation for social reasons: A literature review.
\emph{Healthcare}, 12(23), 2421.

\bibitem[van Loendersloot et~al.(2011)]{loendersloot2011}
van Loendersloot, L. L., Moolenaar, L. M., Mol, B. W. J., Repping, S.,
van der Veen, F., \& Goddijn, M. (2011). Expanding reproductive
lifespan: A cost-effectiveness study on oocyte freezing. \emph{Human
Reproduction}, 26(11), 3054--3060. DOI:
\nolinkurl{10.1093/humrep/der284}.

\bibitem[Luke et~al.(2012)]{luke2012}
Luke, B., Brown, M. B., Wantman, E., Lederman, A., Gibbons, W.,
Schattman, G. L., Lobo, R. A., Leach, R. E., \& Stern, J. E. (2012).
Cumulative birth rates with linked assisted reproductive technology
cycles. \emph{The New England Journal of Medicine}, 366(26), 2483--2491.
DOI: \nolinkurl{10.1056/NEJMoa1110238}.

\bibitem[Matorras and Valladolid(2001)]{matorras2001}
Matorras, R., \& Valladolid, A. (2001). \emph{Costes de la
reproducci\'on asistida en el sistema p\'ublico de salud: experiencia en
el Hospital de Cruces}.

\bibitem[Matorras et~al.(2016)]{matorras2016}
Matorras, R., Villoro, R., Gonz\'alez-Dom\'inguez, A.,
P\'erez-Camarero, S., Hidalgo-Vera, A., \& Polanco, C. (2016).
Long-term fiscal implications of funding assisted reproduction: A
generational accounting model for Spain. \emph{Reproductive
BioMedicine \& Society Online}, 1, 113--122. DOI:
\nolinkurl{10.1016/j.rbms.2016.04.001}.

\bibitem[Matorras Weinig(2011)]{matorrasweinig2011}
Matorras Weinig, R. (2011). \emph{Libro blanco sociosanitario: la
infertilidad en Espa\~na. Situaci\'on actual y perspectivas}. Madrid.

\bibitem[MESS(2006)]{mess2006}
Ministry of Employment and Social Security. (2006). \emph{Informe sobre
sostenibilidad demogr\'afica y gasto social}. Madrid.

\bibitem[Ministerio de Sanidad(2024)]{minsan2024}
Ministerio de Sanidad. (2024). \emph{Estad\'istica de gasto sanitario
p\'ublico}. Madrid.

\bibitem[Moolenaar et~al.(2014)]{moolenaar2014}
Moolenaar, L. M., Connolly, M. P., Huisman, B., Postma, M. J., Hompes,
P. G. A., van der Veen, F., \& Mol, B. W. J. (2014). Costs and
benefits of individuals conceived after IVF: a net tax evaluation in
The Netherlands. \emph{Reproductive BioMedicine Online}, 28(2),
239--245. DOI: \nolinkurl{10.1016/j.rbmo.2013.10.002}.

\bibitem[National Research Council(2012)]{nrc2012}
National Research Council. (2012). \emph{Aging and the macroeconomy:
long-term implications of an older population}. Washington, DC:
National Academies Press.

\bibitem[Navarro Espigares et~al.(2006)]{navarro2006}
Navarro Espigares, J. L., Mart\'inez Navarro, L., Castilla Alcal\'a,
J. A., \& Hern\'andez Torres, E. (2006). Coste de las t\'ecnicas de
reproducci\'on asistida en un hospital p\'ublico. \emph{Gaceta
Sanitaria}, 20(5), 382--391. DOI:
\nolinkurl{10.1157/13093207}.

\bibitem[Noyes et~al.(2009)]{noyes2009}
Noyes, N., Porcu, E., \& Borini, A. (2009). Over 900 oocyte
cryopreservation babies born with no apparent increase in congenital
anomalies. \emph{Reproductive BioMedicine Online}, 18(6), 769--776.
DOI: \nolinkurl{10.1016/S1472-6483(10)60025-9}.

\bibitem[OECD(2024)]{oecd2024}
OECD. (2024). \emph{Family database and social expenditure indicators}.
Paris.

\bibitem[Pai et~al.(2021)]{pai2021}
Pai, H. D., Baid, R., Palshetkar, N. P., Pai, A., Pai, R. D., \&
Palshetkar, R. (2021). Oocyte cryopreservation -- current scenario and
future perspectives: A narrative review. \emph{Journal of Human
Reproductive Sciences}, 14(4), 340--349. DOI:
\nolinkurl{10.4103/jhrs.jhrs_173_21}.

\bibitem[Partido Popular(2025)]{partidopopular2025}
Partido Popular. (2025). \emph{Ponencia pol\'itica}.

\bibitem[P\'erez-Camarero et~al.(2012)]{perezcamarero2012}
P\'erez-Camarero, S., Hidalgo-Vera, A., \& colleagues. (2012).
\emph{Demographic ageing and public finance sustainability in Spain}.

\bibitem[Rindfuss et~al.(2010)]{rindfuss2010}
Rindfuss, R. R., Guilkey, D. K., Morgan, S. P., \& Kravdal, \O.
(2010). Child-care availability and fertility in Norway.
\emph{Population and Development Review}, 36(4), 725--748. DOI:
\nolinkurl{10.1111/j.1728-4457.2010.00355.x}.

\bibitem[SEF(2001)]{sef2001}
Sociedad Espa\~nola de Fertilidad. (2001). \emph{Registro nacional de
actividad 2001}. Madrid.

\bibitem[SEF(2024)]{sef2024registro}
Sociedad Espa\~nola de Fertilidad. (2024). \emph{National activity
registry of the Spanish Fertility Society}. Madrid.

\bibitem[Seiz et~al.(2023)]{seiz2023}
Seiz, M., Eremenko, T., \& Salazar, L. (2023). Socioeconomic
differences in access to and use of medically assisted reproduction
(MAR) in a context of increasing childlessness. \emph{JRC Working
Papers Series on Social Classes in the Digital Age}.

\bibitem[Shenfield et~al.(2017)]{shenfield2017}
Shenfield, F., de Mouzon, J., Scaravelli, G., Kupka, M., Ferraretti,
A. P., Prados, F. J., \& Goossens, V. (2017). Oocyte and ovarian tissue
cryopreservation in European countries: Statutory background, practice,
storage and use. \emph{Human Reproduction Open}, 2017(1), hox003. DOI:
\nolinkurl{10.1093/hropen/hox003}.

\bibitem[Stock and Watson(2020)]{stockwatson2020}
Stock, J. H., \& Watson, M. W. (2020). \emph{Introduction to
econometrics}. Pearson.

\bibitem[Trokoudes et~al.(2011)]{trokoudes2011}
Trokoudes, K. M., Pavlides, C., \& Zhang, X. (2011). Comparison outcome
of fresh and vitrified donor oocytes in an egg-sharing donation program.
\emph{Fertility and Sterility}, 95(6), 1991--1994.

\bibitem[Vallejo-Torres et~al.(2018)]{vallejotorres2018}
Vallejo-Torres, L., Garc\'ia-Lorenzo, B., \& Serrano-Aguilar, P.
(2018). Estimating a cost-effectiveness threshold for the Spanish NHS.
\emph{Health Economics}, 27(4), 746--761. DOI:
\nolinkurl{10.1002/hec.3633}.

\bibitem[World Bank(2025)]{worldbank2025}
World Bank. (2025). \emph{World Development Indicators: Fertility rate,
total (births per woman)}. Washington, DC. URL:
\url{https://data.worldbank.org/indicator/SP.DYN.TFRT.IN}.

\bibitem[Xunta de Galicia(2024)]{xunta2024}
Presidencia de la Xunta de Galicia. (2024, November 4). Rueda anuncia
que la Xunta aprueba el Plan Gallego de Reproducci\'on Humana Asistida
con una inversi\'on de 7,8 M\texteuro{} para ofrecer la cartera de
servicios m\'as amplia de Espa\~na. Santiago de Compostela, Galicia,
Espa\~na.

\bibitem[Yang et~al.(2022)]{yang2022}
Yang, I.-J., Wu, M.-Y., Chao, K.-H., Wei, S.-Y., Tsai, Y.-Y., Huang,
T.-C., \ldots{} Chen, S.-U. (2022). Usage and cost-effectiveness of
elective oocyte freezing: A retrospective observational study.
\emph{Reproductive Biology and Endocrinology}, 20, 123. DOI:
\nolinkurl{10.1186/s12958-022-00996-1}.

\end{thebibliography}
\endgroup

% =======================================================================
% APPENDICES
% =======================================================================
\clearpage
\appendix

% Reset equation numbering for appendices
\setcounter{equation}{0}
\renewcommand{\theequation}{A.\arabic{equation}}

% -----------------------------------------------------------------------
\section{Data and Estimation}\label{app:data}
% -----------------------------------------------------------------------

\subsection{Dataset description and
preprocessing}\label{app:dataset}%
\footnote{Replication materials for this study, including the code and
data underlying the empirical analysis, are available at
\url{https://github.com/alejandrocorchonfranco/reproeconomics}.}

The empirical analysis is conducted on annual demographic and
macroeconomic data for Spain derived from official sources (INE, World
Bank, Human Fertility Database) and integrated into a single
longitudinal dataset covering 1976 to 2022, providing 47 annual
observations. For the macroeconomic series, World Bank data extend to
2022 and are used without extrapolation. One possible issue with an
annual macroeconomic time series of this length is non-stationarity,
because integrating fertility and macroeconomic series of order one
without cointegration testing risks generating spurious regressions
\citep{grangernewbold1974}. Visual checks suggest that the adjTFR and
female labour force participation series are borderline I(1) and GDP
growth borderline I(0). The ITS specification alleviates some of these
concerns by incorporating a deterministic linear trend, through which
the effect of I(1) can be absorbed when the series is
trend-stationary.

The key demographic variables include the adjusted total fertility rate
(adjTFR) and parity metrics (adjTFR1, adjTFR3plus). Age-specific
fertility rates for women aged 35--39 (ASFR\_(35--39)) and 40--44
(ASFR\_(40--44)) are included to capture delayed motherhood. The data
contain a linear time trend (\texttt{time}), the binary break indicator
(\texttt{intervention}) and the interaction term
(\texttt{post\_intervention\_trend}) for identification of structural
breaks. Macroeconomic covariates include the rate of real GDP growth
per capita (\texttt{GDP\_Growth\_PC}), female labour force
participation (\texttt{Female\_LFP}), and the youth unemployment rate
(\texttt{Juvenile\_Unemployment\_Rate}). Due to data constraints, the
latter two variables are observed for shorter subperiods and included
in sensitivity specifications only. All variables are harmonised to
annual frequency and aligned by calendar year; missing macroeconomic
observations are excluded from specifications requiring those controls.

\subsection{Linear ITS model and results}\label{app:linear}

The magnitude of the post-break structural change is quantified by the
coefficients associated with the break terms in
Equation~\eqref{eq:its-model}, where $\beta_2$ measures the
instantaneous level shift at the structural breakpoint and $\beta_3$
captures the change in the slope of the outcome trajectory after the
break. For any post-break period $t \geq T_0$, the estimated dynamic
deviation from the counterfactual is

\begin{equation}
\widehat{\tau}_t = \widehat{\beta}_2 + \widehat{\beta}_3(t - T_0),
\label{eq:dynamic-deviation}
\end{equation}

\noindent which represents the deviation of the observed outcome from
its counterfactual path implied by the pre-break trend. The
population regression representation and exogeneity assumption are,
respectively,

\begin{equation}
Y_t = X_t^{\prime}\beta + u_t
\label{eq:population-model}
\end{equation}

\begin{equation}
E(u_t \mid X_t) = 0.
\label{eq:exogeneity}
\end{equation}

Under the maintained exogeneity assumption, OLS yields unbiased and
consistent estimators of $\beta$; inference uses HAC standard errors
(lag~=~3) throughout.

The regression output is presented in Table~\ref{tab:its-linear} (see
tables at the end of the manuscript). Note that the time trend, the
structural break level shift, and the post-break slope are
statistically significant for most fertility outcomes, particularly for
aggregate fertility (adjTFR), higher-order births (adjTFR3+), and
ASFRs at older maternal ages. First-birth fertility (adjTFR1) shows
weaker explanatory power and limited statistical significance across
most coefficients, consistent with the interpretation that entry into
motherhood is relatively rigid and less responsive to structural breaks
and macroeconomic fluctuations.

For aggregate fertility (adjTFR), the pre-break time trend is negative
and highly significant ($\hat{\beta}_1 = -66.64$; $p<0.01$), confirming
a pronounced secular decline prior to 2010. The level-shift coefficient
at the structural break is also negative and significant
($\hat{\beta}_2 = -1399.35$; $p<0.01$), indicating a sharp downward
discontinuity at $T_0$. The post-break slope coefficient is positive
and significant ($\hat{\beta}_3 = +45.23$; $p<0.01$), implying a
deceleration of the long-run decline after 2010; however, since the
positive post-break slope only partially offsets the initial level
shift, adjTFR remains persistently below its pre-break counterfactual
path throughout the post-2010 period. The coefficient magnitude of
$-1399.35$ reflects units of the adjTFR series (expressed per 1,000
women aged 15--49), which differs from the per-woman scale of ASFR and
parity-specific series; cross-series magnitude comparisons require
standardisation.

GDP per capita growth is negatively and significantly associated with
adjTFR ($\hat{\beta} = -18.47$; $p<0.01$), while female labour force
participation displays a strong positive association
($\hat{\beta} = +54.18$; $p<0.01$). The youth unemployment rate does
not exhibit a statistically significant effect once other factors are
controlled for.

The asymmetry between age groups central to the calibration of $\alpha$
is confirmed: ASFR at ages 35--39 shows a negative level shift
($\hat{\beta}_2 = -0.057$; $p<0.01$) followed by a positive post-break
slope ($\hat{\beta}_3 = +0.00170$; $p<0.01$), consistent with partial
convergence toward the counterfactual. ASFR at ages 40--44 exhibits a
similar level shift ($\hat{\beta}_2 = -0.041$; $p<0.01$) and positive
post-break slope ($\hat{\beta}_3 = +0.00126$; $p<0.01$), but the slope
correction is insufficient to close the gap with the counterfactual
over the estimation horizon.

\subsection{Logit ITS as a qualitative robustness check}\label{app:logit}

To complement the linear ITS, a logit model is estimated for a binary
indicator of low-fertility regime membership $Y_t^L$, equal to 1 when
adjTFR falls within the low-fertility regime and 0 otherwise. The
specification is

\begin{equation}
\Pr(Y_t^L = 1 \mid X_t) = \frac{\exp\!\left(\alpha_0 + \alpha_1 t +
\alpha_2 D_t + \alpha_3 (t-T_0)D_t + \delta'X_t\right)}{1 +
\exp\!\left(\alpha_0 + \alpha_1 t + \alpha_2 D_t +
\alpha_3(t-T_0)D_t + \delta'X_t\right)}.
\label{eq:logit-appendix}
\end{equation}

An important limitation of this specification is that the dependent
variable is constructed endogenously from the same adjTFR series used
in the linear ITS, creating a risk of circular identification. The
near-perfect AUC (= 1.00) is consistent with quasi-complete separation
driven by this circularity; goodness-of-fit statistics (Nagelkerke
$R^2 = 0.988$) should be interpreted as evidence of near-deterministic
classification within the estimation sample rather than robust
out-of-sample predictive performance. The logit results are accordingly
treated as a qualitative robustness check on the direction of the ITS
finding, confirming a lower probability of observing high-fertility
adjTFR outcomes after the structural break, consistent with the regime
shift identified in the linear specification. They do not constitute an
independent source of identification.

% Reset figure counter for Appendix B
\setcounter{figure}{0}
\renewcommand{\thefigure}{B.\arabic{figure}}

% Reset equation numbering for Appendix B
\setcounter{equation}{0}
\renewcommand{\theequation}{B.\arabic{equation}}

% -----------------------------------------------------------------------
\section{Decision-Tree Representation of the Cost Model}\label{app:decisiontree}
% -----------------------------------------------------------------------

This appendix presents a stylised decision tree of the economic logic
underpinning the publicly funded elective OC programme analysed in the
main text. The model, as described in
Sections~\ref{sec:assumptions} and~\ref{sec:parameters}, combines
behavioural parameters (uptake and utilisation rates) with clinical
outcomes (success rates) to estimate the expected fiscal impact of the
policy.

The figure establishes the sequential structure within which the
framework in Section~\ref{sec:benchmark} operates: eligibility in the
target population, the choice to undergo vitrification, the probability
of using stored oocytes, and the conditional probability of obtaining a
live birth. These are the parameters used to interpret the estimated
cost-effectiveness as presented in
Sections~\ref{sec:cost-per-birth} and~\ref{sec:fiscal}.

Figure~\ref{fig:decision_tree} clarifies how programme costs are
incurred even when oocytes are never ultimately used, and that the
fiscal break-even condition is jointly determined by behavioural and
clinical margins. It should be treated as a schematic example of the
analytic framework in Section~\ref{sec:budget}, not as a standalone
empirical outcome.

% =======================================================================
% TABLES (after bibliographical references)
% =======================================================================
\clearpage

% Reset table counter for Appendix A
\setcounter{table}{0}
\renewcommand{\thetable}{A.\arabic{table}}

\begin{sidewaystable}[p]
\centering
\caption{Linear ITS estimates of fertility outcomes}
\label{tab:its-linear}
\begin{threeparttable}
\footnotesize
\setlength{\tabcolsep}{5pt}
\renewcommand{\arraystretch}{1.28}
\begin{adjustbox}{max width=\textwidth}
\begin{tabular}{lcccccccc}
\toprule
Outcome & $N$ & $R^2$ & Time Trend ($\beta_1$) & Intervention
($\beta_2$) & Post-Intervention Trend ($\beta_3$) & GDP Growth PC &
Female LFP & Youth Unemployment \\
\midrule
adjTFR & 47 & 0.949 &
\makecell[r]{$-66.64^{***}$\\(5.29)} &
\makecell[r]{$-1399.35^{***}$\\(407.50)} &
\makecell[r]{$+45.23^{***}$\\(9.26)} &
\makecell[r]{$-18.47^{***}$\\(6.04)} &
\makecell[r]{$+54.18^{***}$\\(11.61)} &
\makecell[r]{0.82\\(2.74)} \\
adjTFR1 & 47 & 0.539 &
\makecell[r]{$-0.010$\\(0.006)} &
\makecell[r]{0.45\\(0.41)} &
\makecell[r]{$-0.010$\\(0.009)} &
\makecell[r]{$-0.014^{**}$\\(0.006)} &
\makecell[r]{0.006\\(0.013)} &
\makecell[r]{0.000\\(0.003)} \\
adjTFR3+ & 47 & 0.974 &
\makecell[r]{$-0.042^{***}$\\(0.004)} &
\makecell[r]{$-1.55^{***}$\\(0.21)} &
\makecell[r]{$+0.046^{***}$\\(0.005)} &
\makecell[r]{$-0.002$\\(0.002)} &
\makecell[r]{$+0.039^{***}$\\(0.005)} &
\makecell[r]{$-0.002$\\(0.001)} \\
ASFR 35--39 & 47 & 0.949 &
\makecell[r]{$-0.00177^{***}$\\(0.00031)} &
\makecell[r]{$-0.057^{***}$\\(0.018)} &
\makecell[r]{$+0.00170^{***}$\\(0.00044)} &
\makecell[r]{$-0.00008$\\(0.00015)} &
\makecell[r]{$+0.00385^{***}$\\(0.00047)} &
\makecell[r]{$-0.00025^{**}$\\(0.00010)} \\
ASFR 40--44 & 47 & 0.970 &
\makecell[r]{$-0.00090^{***}$\\(0.00008)} &
\makecell[r]{$-0.041^{***}$\\(0.005)} &
\makecell[r]{$+0.00126^{***}$\\(0.00013)} &
\makecell[r]{$-0.00003$\\(0.00005)} &
\makecell[r]{$+0.00128^{***}$\\(0.00011)} &
\makecell[r]{$-0.00005^{**}$\\(0.00003)} \\
\bottomrule
\end{tabular}
\end{adjustbox}
\begin{tablenotes}[flushleft]
\footnotesize
\item[a] HAC standard errors (lag = 3) are reported in parentheses.
$^{***}p<0.01$, $^{**}p<0.05$, $^{*}p<0.10$.
\item[b] The larger magnitude of the adjTFR coefficients reflects a
difference in scale: adjTFR is expressed per 1,000 women aged 15--49,
whereas adjTFR1 and adjTFR3+ are expressed as rates per woman.
Cross-row comparisons of coefficient magnitude are therefore not
meaningful without standardisation.
\end{tablenotes}
\end{threeparttable}
\end{sidewaystable}

% =======================================================================
% FIGURES (after bibliographical references and tables)
% =======================================================================
\clearpage

% Reset figure counter for main-text figures
\setcounter{figure}{0}
\renewcommand{\thefigure}{\arabic{figure}}

\begin{figure}[H]
\centering
\includegraphics[width=0.88\textwidth]{media/media/image3.png}
\caption{Pre-trend test for adjusted total fertility: observed versus
predicted pre-intervention trajectory}
\label{fig:pretrend}
\caption*{\footnotesize\textit{Note:} The solid line shows the
observed trajectory of adjTFR and the dashed line shows the estimated
pre-break linear trend using only pre-2010 data. Close alignment
supports the counterfactual assumption.\\
\textit{Source:} Author's calculations based on INE and Human
Fertility Database (1976--2022).}
\end{figure}

\clearpage

\begin{figure}[H]
\centering
\includegraphics[width=0.88\textwidth]{media/media/image2.png}
\caption{Sensitivity of estimated structural-break coefficient to
alternative candidate break-dates in the ITS model}
\label{fig:break_sensitivity}
\caption*{\footnotesize\textit{Note:} The figure reports the estimated
level-shift coefficient ($\beta_2$) from the ITS specification under
alternative candidate break-dates. Coefficients are plotted with
corresponding confidence intervals.\\
\textit{Source:} Author's calculations.}
\end{figure}

\clearpage

\begin{figure}[H]
\centering
\includegraphics[width=0.88\textwidth]{media/media/image4.png}
\caption{Post-intervention gap in ASFRs, controlling for
macroeconomic variables}
\label{fig:asfr_gap}
\caption*{\footnotesize\textit{Note:} The figure illustrates the
estimated post-break gap in age-specific fertility rates for women
aged 35--39 and 40--44, derived from the ITS specification including
macroeconomic controls. The horizontal dashed line represents the null
deviation (gap = 0), while the series capture the cumulative
divergence from the pre-break counterfactual trajectory after $T_0 =
2010$.\\
\textit{Source:} Author's calculations based on INE and Human
Fertility Database (1976--2022).}
\end{figure}

\clearpage

\begin{figure}[H]
\centering
\includegraphics[width=0.88\textwidth]{media/media/image9.png}
\caption{Cost per additional birth and sensitivity to utilisation rate
and success rate}
\label{fig:cost_sensitivity}
\caption*{\footnotesize\textit{Note:} Bidimensional sensitivity
analysis of the cost per additional live birth over the plausible
ranges of utilisation rate ($u$) and success rate per thaw cycle
($s$), with demographic additionality parameter $\alpha = 0.50$.
Values expressed in 2023 euros.\\
\textit{Source:} Author's calculations.}
\end{figure}

\clearpage

\begin{figure}[H]
\centering
\includegraphics[width=0.88\textwidth]{media/media/image10.png}
\caption{Deterministic sensitivity analysis of cost per additional
birth}
\label{fig:tornado}
\caption*{\footnotesize\textit{Note:} Each horizontal bar shows the
range of the cost per additional live birth obtained by varying one
parameter at a time over its indicated interval while holding all
other parameters at their base-case values. The dashed black vertical
line denotes the base-case estimate (\euro{}433,333); the dotted blue
vertical line denotes the NFC break-even threshold (\euro{}91,391).
All values in 2023 euros.\\
\textit{Source:} Author's calculations.}
\end{figure}

\clearpage

\begin{figure}[H]
\centering
\includegraphics[width=0.88\textwidth]{media/media/image11.png}
\caption{Annual budgetary commitments under alternative uptake
scenarios}
\label{fig:budget_scenarios}
\caption*{\footnotesize\textit{Note:} Annual programme costs (2023
euros) under 5\%, 10\%, and 20\% uptake assumptions, including
vitrification, thaw-and-transfer IVF, and administration costs.
Adverse event costs are excluded.\\
\textit{Source:} Author's calculations based on SEF (2024) and
Ministerio de Sanidad (2024).}
\end{figure}

\clearpage

\begin{figure}[H]
\centering
\includegraphics[width=0.88\textwidth]{media/media/image12.png}
\caption{Estimated OC utilisation rates by income quintile under
alternative funding designs}
\label{fig:quintiles}
\caption*{\footnotesize\textit{Note:} OC utilisation rates by income
quintile under three designs: the current private market, a universal
public subsidy, and a means-tested subsidy (Q1--Q3). Calibrated on the
access gradients reported by \citet{alonpinilla2021} and
\citet{seiz2023}.\\
\textit{Source:} Author's calculations based on Alon and Pinilla
(2021) and Seiz et al. (2023).}
\end{figure}

\clearpage

\begin{figure}[H]
\centering
\includegraphics[width=0.88\textwidth]{media/media/image13.png}
\caption{Concentration curves for programme distributive analysis}
\label{fig:concentration}
\caption*{\footnotesize\textit{Note:} Concentration curves for the
three programme designs. The dashed 45-degree line represents perfect
equality. A positive concentration index (CI) indicates a pro-rich
distribution. The calibration is anchored to Spain's income
distribution (World Bank, WDI 2023) and to the income gradient in
private access reported by \citet{alonpinilla2021}: private market
CI~$\approx$~0.42; universal subsidy CI~$\approx$~0.18; means-tested
subsidy CI~$\approx$~0.05.\\
\textit{Source:} Author's calculations based on Alon and Pinilla
(2021), Seiz et al. (2023), and World Bank WDI (2023).}
\end{figure}

\clearpage

% Appendix B figure -- reset numbering to B.1
\setcounter{figure}{0}
\renewcommand{\thefigure}{B.\arabic{figure}}

\begin{figure}[H]
\centering
\includegraphics[width=0.88\textwidth]{media/media/image8.png}
\caption{Decision-tree representation of the cost structure of a
publicly funded elective OC programme}
\label{fig:decision_tree}
\caption*{\footnotesize\textit{Note:} A simplified decision tree
modelling the expected fiscal cost per live birth under a publicly
funded elective OC programme. The baseline scenario assumes an uptake
rate of 8\%, a utilisation rate of 8\%, and a clinical success rate of
30\%. Costs include the vitrification procedure (approximately
\euro{}5,000 per patient) and the thaw-and-transfer IVF cycle
(approximately \euro{}2,500). NFC refers to the net fiscal contribution
per additional birth estimated in the main text. The fiscal break-even
condition is achieved when the expected cost per live birth does not
exceed the NFC (see Section~\ref{sec:fiscal}).\\
\textit{Source:} Author's own elaboration.}
\end{figure}

\end{document}
